\clearpage
\FrontTitle{KATA PENGANTAR}
\addcontentsline{toc}{chapter}{KATA PENGANTAR}
Laporan ini disusun sebagai dokumentasi pelaksanaan kerja praktik di PT Akhdani Reka Solusi pada \PeriodeKP. Selama periode tersebut, penulis memperoleh kesempatan untuk berkontribusi pada pekerjaan pengembangan perangkat lunak serta memahami proses kerja dalam lingkungan profesional.

Pelaksanaan kerja praktik menjadi sarana untuk menerapkan pengetahuan informatika pada kebutuhan nyata. Selain keterampilan teknis, pekerjaan menuntut kemampuan memahami penugasan, berkomunikasi dengan rekan kerja, dan menyesuaikan diri terhadap perubahan. Pengalaman serta hasil kegiatan tersebut dituangkan dalam laporan ini sebagai bahan evaluasi dan pembelajaran.

Penyusunan laporan dan pelaksanaan kerja praktik didukung oleh berbagai pihak. Untuk itu, penulis menyampaikan terima kasih kepada:
\begin{enumerate}[label=\arabic*)]
\item \NamaDosenPembimbing{} selaku dosen pembimbing kerja praktik;
\item \NamaPembimbingPerusahaan{} selaku pembimbing perusahaan dan \JabatanPembimbingPerusahaan;
\item \NamaPembimbingTeknis{} selaku pembimbing teknis dengan peran \JabatanPembimbingTeknis;
\item rekan-rekan di PT Akhdani Reka Solusi yang membantu pelaksanaan pekerjaan; serta
\item Program Studi Teknik Informatika dan staf administrasi Institut Teknologi Bandung yang mendukung pelaksanaan kerja praktik.
\end{enumerate}
Penulis berharap laporan ini dapat menjelaskan pelaksanaan kerja praktik dengan baik dan menjadi bahan perbaikan bagi pengembangan kemampuan profesional berikutnya.

\vspace{0.7cm}
\begin{flushright}
Bandung, \Isian{tanggal penyelesaian laporan}
\end{flushright}
\vspace{0.5cm}
\begin{center}
\begin{tabular}{cc}
Penulis I & Penulis II\\[1.4cm]
\NamaRaif & \NamaFrederiko\\
\NIMRaif & \NIMFrederiko
\end{tabular}
\end{center}
